\documentclass[11pt,a4paper,reqno]{amsart}
\usepackage[T1]{fontenc}
\usepackage{lmodern}
\usepackage[a4paper,margin=28mm]{geometry}
\usepackage{amsmath,amssymb,amsthm,mathtools,microtype}
\usepackage[hidelinks,unicode]{hyperref}
\hypersetup{pdftitle={Pattern complexity and Nivat's conjecture},pdfsubject={Rectangular pattern complexity and polynomial annihilators},pdfauthor={GPT-6 Pro}}
\newtheorem{theorem}{Theorem}[section]
\newtheorem{proposition}[theorem]{Proposition}
\newtheorem{lemma}[theorem]{Lemma}
\newtheorem{corollary}[theorem]{Corollary}
\theoremstyle{remark}
\newtheorem{remark}[theorem]{Remark}
\newtheorem{example}[theorem]{Example}
\newcommand{\Z}{\mathbb Z}
\newcommand{\Q}{\mathbb Q}
\newcommand{\R}{\mathbb R}
\newcommand{\LL}{\mathcal L}
\newcommand{\mon}[1]{\boldsymbol X^{#1}}
\newcommand{\De}{\Delta}
\newcommand{\restr}[2]{\left.#1\right|_{#2}}
\DeclareMathOperator{\Ann}{Ann}
\DeclareMathOperator{\supp}{supp}
\DeclareMathOperator{\Pat}{Pat}
\DeclareMathOperator{\spanQ}{span_{\Q}}
\numberwithin{equation}{section}
\setlength{\emergencystretch}{2em}
\allowdisplaybreaks[1]
\title[Pattern complexity and Nivat's conjecture]{Pattern complexity and Nivat's conjecture}
\author{GPT-6 Pro}
\date{September 14, 2026}
\subjclass[2020]{37B10, 68R15}
\keywords{Nivat's conjecture, rectangular complexity, polynomial annihilators, symbolic dynamics}
\begin{document}
\begin{abstract}
We prove Nivat's conjecture: a configuration on $\Z^2$ over a finite alphabet with at most $mn$ distinct $m\times n$ patterns, for some positive integers $m,n$, has a nonzero global period. After labeling the alphabet by rationals, we construct a Laurent polynomial operator that reduces the pattern count by at least the decrease in the window's area. The construction identifies a line polynomial generating the full annihilator ideal of the difference of two distinct orbit-closure points agreeing on a half-plane. Induction on rectangle area then reduces the conjecture to configurations annihilated by a product of two difference operators. We prove this case using the boundary-extension methods of Cyr--Kra and Szabados. The full proof is formalized in Lean using mathlib.
\end{abstract}
\maketitle

\section{Introduction}\label{sec:intro}
Let $A$ be a finite alphabet and let $c:\Z^2\to A$ be a configuration. For a finite set $D\subset\Z^2$, its set of occurring patterns is
\[
 \Pat_c(D)=\{(c(z+u))_{u\in D}:z\in\Z^2\},
 \qquad P_c(D)=|\Pat_c(D)|.
\]
Write
\[
 R_{m,n}=\{0,\ldots,m-1\}\times\{0,\ldots,n-1\},
 \qquad P_c(m,n)=P_c(R_{m,n}).
\]
We call $c$ \emph{periodic} if there is an $h\in\Z^2\setminus\{0\}$ such that $c(z+h)=c(z)$ for every $z$. Only one period vector is required.

The Morse--Hedlund theorem says that a bi-infinite word is periodic if and only if, for some $n$, it has at most $n$ distinct words of length $n$~\cite{MH}. Nivat's conjecture is the analogous implication for rectangles in two dimensions.

\begin{theorem}\label{thm:main}
Let $c:\Z^2\to A$ be a configuration over a finite alphabet. If $P_c(m,n)\le mn$ for some positive integers $m,n$, then $c$ is periodic.
\end{theorem}

Our proof uses two consequences of a low pattern count. After a rational labeling of the alphabet, the occurring patterns satisfy an affine relation; taking a difference gives a polynomial annihilator. This is the starting point of the algebraic approach of Kari and Szabados~\cite{KS}. There is also a combinatorial consequence: if adjoining a boundary adds too few patterns, some boundary values are uniquely determined. This observation underlies the generating-set and balanced-set arguments of Cyr and Kra~\cite{CK}. Szabados used these methods to prove Nivat's conjecture for sums of two periodic configurations~\cite{Sz}.

Kari and Szabados already relate linear dependence of patterns to polynomial annihilators~\cite[proofs of Lemmas~6 and~32]{KS}. Here we use this relation for differences of patterns with the same polynomial image. The resulting estimate accounts for every coordinate lost when the polynomial is applied to a rectangle, and makes a reduction in area possible. We describe it after fixing notation. Section~\ref{sec:twofactor} then gives the two-periodic-summand argument in the form needed to complete the reduction.

\subsection{Notation}
An injective labeling of $A$ by rational numbers preserves pattern counts and periods. Except where stated otherwise, configurations below have finite range in $\Q$. This class is closed under rational linear combinations and translations.

For $u\in\Z^2$, let
\[
 (T^u c)(z)=c(z+u),\qquad \De_u=T^u-I.
\]
These shift and difference operators commute. Set
\[
 \LL=\Q[X^{\pm1},Y^{\pm1}],\qquad
 \mon u=X^{u_1}Y^{u_2}.
\]
A Laurent polynomial $f=\sum_u f_u\mon u$ acts by
\[
 (f(T)c)(z)=\sum_u f_uc(z+u).
\]
The sum is finite. We write $\supp f=\{u:f_u\ne0\}$ and
\[
 \Ann(c)=\{f\in\LL:f(T)c=0\}
\]
for the support of $f$ and the annihilator ideal of $c$, respectively. Thus $h\ne0$ is a period of $c$ exactly when $\mon h-1\in\Ann(c)$. Our use of forward shifts fixes the sign convention for the polynomial action.

A nonzero vector $v\in\Z^2$ is \emph{primitive} if $\gcd(|v_1|,|v_2|)=1$. Such a vector extends to a basis $v,w$ of $\Z^2$. Following~\cite{KS}, a \emph{line polynomial} is a Laurent polynomial whose support contains at least two points and lies on a line. Up to a monomial factor, it has the form $\phi(\mon v)$ for a primitive $v$ and a nonconstant $\phi\in\Q[Z]$ with $\phi(0)\ne0$.

For later counting arguments, put
\[
 \delta_c(D)=P_c(D)-|D|.
\]
This is the discrepancy of the window $D$ in the terminology of~\cite[Section~2.2]{CK}. There is one pattern on the empty set, so $P_c(\varnothing)=1$. Translating a window does not change its complexity. Restriction from a larger window to a smaller one is a surjection between the corresponding sets of occurring patterns.

Let $X_c$ be the closure of $\{T^u c:u\in\Z^2\}$ in the product space $A^{\Z^2}$, where $A=c(\Z^2)$ has the discrete topology. Every finite pattern of a point of $X_c$ occurs in $c$. Moreover,
\begin{equation}\label{eq:inheritance}
 f\in\Ann(c),\quad x,y\in X_c
 \quad\Longrightarrow\quad f(T)x=f(T)y=f(T)(x-y)=0.
\end{equation}
Both assertions follow by taking limits on finitely many coordinates. The space $X_c$ is compact; sequential compactness can be proved by diagonal extraction over the countable lattice.

\subsection{The complexity inequality and the reduction in area}
Fix $0\ne\Phi\in\LL$ and a rectangle $R$, and let
\begin{equation}\label{eq:erosion-definition}
 S=\{z\in\Z^2:z+\supp\Phi\subseteq R\}.
\end{equation}
A pattern on $R$ determines the restriction of $\Phi(T)c$ to $S$. Therefore
$P_{\Phi(T)c}(S)\le P_c(R)$. This inequality alone is insufficient: when $|S|<|R|$, it need not imply $P_{\Phi(T)c}(S)\le|S|$.

We prove the stronger estimate
\begin{equation}\label{eq:intro-descent}
 P_{\Phi(T)c}(S)\le P_c(R)-|R|+|S|
\end{equation}
under the hypothesis that there are $x,y\in X_c$ with
\[
 d=x-y\ne0,\qquad \Ann(d)=(\Phi).
\]
Here $(\Phi)$ denotes the principal ideal generated by $\Phi$. To see why the hypothesis is relevant, let $F:\Q^R\to\Q^S$ be the map induced by $\Phi(T)$. Every translated restriction of $d$ is the difference of two occurring $R$-patterns with the same image under $F$. The exact ideal equality implies that these translated restrictions span \emph{all} of $\ker F$, whose dimension is $|R|-|S|$. A fiber containing $k$ patterns supplies at most $k-1$ generators for its differences. Summing over the fibers gives~\eqref{eq:intro-descent}.

The rest of the argument constructs a line polynomial for which this ideal equality holds. A compactness argument gives distinct $x,y\in X_c$ agreeing on a half-plane. A theorem of Kari--Szabados provides a product of difference operators annihilating $d=x-y$; the form needed here is proved in Appendix~\ref{app:product}. On configurations vanishing sufficiently far into the half-plane, difference operators transverse to its boundary are injective. Cancelling those factors shows that $d$ has a period $qv$ parallel to the boundary. Section~\ref{sec:ideal} then proves
\[
 \Ann(d)=\bigl(\phi(\mon v)\bigr),\qquad \phi\mid Z^q-1,
\]
where $\phi$ is the monic polynomial of least degree satisfying $\phi(T^v)d=0$. The half-plane condition rules out additional annihilators involving the transverse shift.

Choose a counterexample with a witnessing rectangle of least area, allowing all finite rational alphabets. For $\Phi=\phi(\mon v)$, the set $S$ in~\eqref{eq:erosion-definition} is a smaller rectangle. Inequality~\eqref{eq:intro-descent} makes it a low-complexity window for $c'=\Phi(T)c$. It cannot be empty, since that would give $1\le0$. Minimality therefore makes $c'$ periodic. If $t\ne0$ is one of its periods, the divisibility $\phi\mid Z^q-1$ gives
\[
 \De_t\De_{qv}c=0.
\]
Section~\ref{sec:twofactor} proves that this identity, together with the original rectangular complexity bound, forces $c$ to be periodic. Section~\ref{sec:completion} assembles these steps into the proof of Theorem~\ref{thm:main}.

\section{Pattern complexity under a polynomial operator}\label{sec:descent}
We first prove~\eqref{eq:intro-descent}. For a finite set $D$, identify $\Q^D$ with the Laurent polynomials supported in $D$ and use the pairing $\langle p,g\rangle_D=\sum_{z\in D}p_zg_z$. The rectangle assumption enters in determining which multiples of $\Phi$ can be supported in $R$.

\begin{lemma}\label{lem:supported}
Let $R$ be a finite axis-aligned rectangle, let $0\ne\Phi\in\LL$, and set
$S=\{z:z+\supp\Phi\subseteq R\}$. Then
\begin{equation}\label{eq:supported}
 (\Phi)\cap\Q^R=\Phi\Q^S,
 \qquad \dim_\Q\bigl((\Phi)\cap\Q^R\bigr)=|S|.
\end{equation}
\end{lemma}
\begin{proof}
Multiplication by $\Phi$ sends $\Q^S$ into $\Q^R$. Conversely, suppose that $\Phi g$ is nonzero and supported in $R$. The minimum and maximum exponents in each coordinate add under multiplication. For example, regarding the factors as Laurent polynomials in $X$ over the integral domain $\Q[Y^{\pm1}]$ shows that their leading coefficients have nonzero product. Writing $\max_X f$ for the largest $X$-exponent in $\supp f$, we obtain
\[
 \max_X(\Phi g)=\max_X\Phi+\max_Xg.
\]
The same holds for the minimum exponent and for the $Y$ coordinate.

Write the coordinate intervals of $R$ as $[a_j,b_j]\cap\Z$, and let $\alpha_j,\beta_j$ be the minimum and maximum $j$-coordinates of $\supp\Phi$. Every $z\in\supp g$ satisfies
\[
 a_j-\alpha_j\le z_j\le b_j-\beta_j\qquad(j=1,2).
\]
For a rectangle, these inequalities are exactly $z+\supp\Phi\subseteq R$. Hence $g\in\Q^S$. Multiplication by $\Phi$ is injective, so its image has dimension $|S|$. The same statements hold when $S$ is empty.
\end{proof}

\begin{theorem}\label{thm:descent}
Let $c$ have finite rational range, and suppose that $x,y\in X_c$ satisfy
\[
 d=x-y\ne0,\qquad \Ann(d)=(\Phi)
\]
for some nonzero $\Phi\in\LL$. For a nonempty rectangle $R$, set
$S=\{z:z+\supp\Phi\subseteq R\}$. Then
\begin{equation}\label{eq:descent}
 P_{\Phi(T)c}(S)-|S|\le P_c(R)-|R|.
\end{equation}
\end{theorem}
\begin{proof}
Write $\Phi=\sum_u\Phi_u\mon u$. The induced linear map is
\[
 F:\Q^R\longrightarrow\Q^S,\qquad
 F(p)(z)=\sum_u\Phi_u p(z+u).
\]
Its transpose is multiplication by $\Phi$: for $p\in\Q^R$ and $g\in\Q^S$,
\[
 \langle Fp,g\rangle_S
 =\sum_{z\in S}\sum_u g_z\Phi_u p(z+u)
 =\langle p,\Phi g\rangle_R.
\]
Consequently $F^{\mathsf T}$ is injective and
\begin{equation}\label{eq:filter-dimension}
 \dim\ker F=|R|-|S|.
\end{equation}

Let
\[
 V_R(d)=\spanQ\{\restr{T^u d}{R}:u\in\Z^2\}.
\]
A vector $b=(b_z)_{z\in R}$ is perpendicular to this space exactly when
\[
 \sum_{z\in R}b_zd(u+z)=0\qquad\text{for every }u\in\Z^2,
\]
that is, when its associated polynomial belongs to $\Ann(d)$. Lemma~\ref{lem:supported} therefore gives
\[
 V_R(d)^\perp=\Ann(d)\cap\Q^R
 =(\Phi)\cap\Q^R=\Phi\Q^S=\operatorname{im}F^{\mathsf T}.
\]
Taking orthogonal complements yields
\begin{equation}\label{eq:full-kernel}
 V_R(d)=\ker F.
\end{equation}

It remains to compare this dimension with the number of occurring patterns. Put $\mathcal P=\Pat_c(R)$. Since polynomial operators commute with translations,
\[
 F(\mathcal P)=\Pat_{\Phi(T)c}(S).
\]
In each fiber choose a representative, and subtract it from every other pattern in that fiber. The resulting list has
\[
 |\mathcal P|-|F(\mathcal P)|=P_c(R)-P_{\Phi(T)c}(S)
\]
vectors and spans every difference between patterns in the same fiber.

For each $u$, the patterns $\restr{T^u x}{R}$ and $\restr{T^u y}{R}$ occur in $c$. They have the same image under $F$, since $\Phi(T)d=0$. The span of the preceding list thus contains $V_R(d)=\ker F$. By~\eqref{eq:filter-dimension},
\[
 |R|-|S|\le P_c(R)-P_{\Phi(T)c}(S),
\]
as required.
\end{proof}

For $\Phi=X-1$, the map $F$ takes horizontal differences on an $m\times n$ rectangle, and its kernel consists of the patterns constant along each row. Its dimension is $n$, the number of sites lost on passing to an $(m-1)\times n$ rectangle. The proof works because the translated restrictions of $x-y$ span all of these row-constant patterns. Knowing only $\Phi(T)(x-y)=0$ would give an inclusion in the kernel, not the equality needed to count $n$ independent directions.

\begin{corollary}\label{cor:line-descent}
In Theorem~\ref{thm:descent}, suppose
\[
 \Phi=\phi(\mon v),\qquad
 \phi(Z)=a_0+a_1Z+\cdots+a_\ell Z^\ell,
 \qquad a_0a_\ell\ne0,
\]
where $v\ne0$ and $\ell\ge1$. If $R=R_{m,n}$ and $P_c(R)\le mn$, then $S$ is a nonempty rectangle of smaller area and $P_{\Phi(T)c}(S)\le|S|$.
\end{corollary}
\begin{proof}
The support of $\Phi$ contains $0$ and $\ell v$ and lies on the segment between them. Since $R$ is a rectangle,
\[
 S=R\cap(R-\ell v),\qquad
 |S|=\max(0,m-\ell|v_1|)\max(0,n-\ell|v_2|).
\]
Theorem~\ref{thm:descent} gives $P_{\Phi(T)c}(S)\le|S|$. If $S$ were empty, this would say $1\le0$. At least one side length strictly decreases because $v\ne0$ and $\ell\ge1$.
\end{proof}

\section{Agreement on a half-plane}\label{sec:halfplane}
We seek two points of $X_c$ to which Theorem~\ref{thm:descent} can be applied. First we construct distinct points that agree on a half-plane. The product-annihilator theorem then shows that their difference is periodic parallel to the boundary. Neither point need be periodic individually.

\subsection{A compactness argument}
The existence of half-plane agreement is a symbolic consequence of the results of Boyle and Lind~\cite[Theorem~3.7]{BL} and Einsiedler, Lind, Miles, and Ward~\cite[Lemma~2.9]{ELMW}. We give a direct proof that also places a disagreement at the origin.

\begin{lemma}\label{lem:halfplane-pair}
If a finite-alphabet configuration $c$ has infinitely many distinct translates, there are $x,y\in X_c$ and a unit vector $\nu\in\R^2$ such that
\begin{equation}\label{eq:halfplane-pair}
 x(0)\ne y(0),\qquad x(z)=y(z)\quad\text{when }\nu\cdot z<0.
\end{equation}
\end{lemma}
\begin{proof}
For each positive integer $j$, choose distinct translates $x_j,y_j$ agreeing on the lattice points in the Euclidean ball of radius $j$. This is possible because there are only finitely many patterns on that ball. Let $p_j$ be a nearest disagreement and put $r_j=\|p_j\|>j$. A nearest disagreement exists because squared lattice norms are nonnegative integers.

By diagonal extraction and compactness of the unit circle, pass to a subsequence such that
\[
 T^{p_j}x_j\longrightarrow x,\qquad
 T^{p_j}y_j\longrightarrow y,\qquad
 p_j/r_j\longrightarrow\nu,\quad \|\nu\|=1.
\]
Since the alphabet is finite, we may also fix the unequal pair of symbols at the shifted origin. Thus $x(0)\ne y(0)$.

Fix $z$ with $\nu\cdot z<0$. For all sufficiently large $j$,
\[
 \|p_j+z\|^2-r_j^2
 =2r_j\bigl((p_j/r_j)\cdot z\bigr)+\|z\|^2<0.
\]
By the choice of $p_j$, the configurations $x_j,y_j$ agree at $p_j+z$. Passing to the limit gives $x(z)=y(z)$.
\end{proof}

Geometrically, the translated agreement balls approach a half-plane with boundary through the fixed disagreement at the origin. The argument gives no agreement on that boundary; nor does it yet imply that the boundary has rational slope.

\subsection{Polynomial annihilators}
We recall the elementary implication from low complexity to an annihilator~\cite[Lemma~5]{KS}.

\begin{lemma}\label{lem:ann-exists}
If $c$ is rational-valued and $P_c(D)\le|D|$ for a nonempty finite set $D$, then $\Ann(c)\ne\{0\}$.
\end{lemma}
\begin{proof}
Put $N=|D|$. The at most $N$ augmented vectors $(1,p)\in\Q^{N+1}$, for $p\in\Pat_c(D)$, have a nonzero common orthogonal vector $(\beta,(b_u)_{u\in D})$. Some $b_u$ is nonzero, since otherwise $\beta$ would also vanish. Hence
\[
 g=\sum_{u\in D}b_u\mon u\ne0,
 \qquad g(T)c=-\beta\mathbf1,
\]
where $\mathbf1$ is the constant configuration with value one. For any $h\ne0$, the polynomial $(\mon h-1)g$ is nonzero and annihilates $c$.
\end{proof}

We use the following consequence of the Kari--Szabados product-annihilator theorem~\cite[Theorem~11 and Corollary~12]{KS}. Appendix~\ref{app:product} gives a proof of this version from their prime-dilation argument. Unlike the full statement of their corollary, it allows repeated directions.

\begin{proposition}[Kari--Szabados]\label{prop:product}
If a nonzero finite-range rational configuration $c$ has a nonzero annihilator, then
\begin{equation}\label{eq:product-ann}
 \De_{h_1}\cdots\De_{h_s}c=0
\end{equation}
for some $s\ge1$ and nonzero $h_1,\ldots,h_s\in\Z^2$.
\end{proposition}

For finite-range configurations, the removal of repeated factors below is a special case of~\cite[Lemmas~9--10]{KS}. The following bounded-sequence argument is all that we need.

\begin{lemma}\label{lem:repeated}
Let $e$ be a bounded rational configuration, let $h\ne0$, and let $s\ge1$. If $\De_h^s e=0$, then $\De_h e=0$.
\end{lemma}
\begin{proof}
For each $z$, set $a_j=e(z+jh)$. If the $s$th difference of this bounded sequence is zero, its $(s-1)$st difference is constant. For $s\ge2$, a nonzero such constant would make the bounded $(s-2)$nd difference grow linearly. The constant is therefore zero. Repetition proves that the first difference vanishes.
\end{proof}

\subsection{A period parallel to the boundary}
We use the half-plane cancellation argument of Kari and Moutot~\cite[Proposition~13 and the proof of Lemma~17]{KM}. It removes every difference factor transverse to the boundary; boundedness then deals with any repeated factors in the remaining direction.

\begin{proposition}\label{prop:tangent-period}
Suppose $d\ne0$ has finite rational range and a nonzero annihilator. If $d(z)=0$ whenever $\nu\cdot z<0$, where $\nu\ne0$, then there are a primitive $v\in\Z^2$ and $q\ge1$ such that
\[
 \nu\cdot v=0,\qquad T^{qv}d=d.
\]
\end{proposition}
\begin{proof}
Let $\mathcal V_\nu$ be the vector space of rational configurations $e$ for which, for some $b\in\R$,
\[
 e(z)=0\quad\text{whenever }\nu\cdot z<b.
\]
Allowing $b$ to depend on $e$ makes this space invariant under every shift. If $\nu\cdot h\ne0$, then $\De_h$ is injective on $\mathcal V_\nu$: an $h$-periodic configuration in this space vanishes, since every $h$-orbit enters its vanishing half-plane.

Apply Proposition~\ref{prop:product} to $d$. The difference operators commute and preserve $\mathcal V_\nu$, so each factor transverse to $\nu^\perp$ can be cancelled by this injectivity. Some factor remains, since otherwise $d=0$. All remaining vectors lie on the line $\nu^\perp$. Choose a primitive lattice vector $v$ along it and write the resulting relation as
\[
 \prod_{i=1}^r(T^{q_i v}-I)d=0,
 \qquad q_i\in\Z\setminus\{0\}.
\]
Let $q$ be a positive common multiple of the $|q_i|$. In $\Q[Z^{\pm1}]$, each $Z^{q_i}-1$ divides $Z^q-1$. Their product therefore divides $(Z^q-1)^r$, and $\De_{qv}^{\,r}d=0$. Lemma~\ref{lem:repeated} gives $\De_{qv}d=0$.
\end{proof}

Cancellation here takes place on $\mathcal V_\nu$, not on the space of all configurations, where a difference operator is certainly not injective. The argument is also specifically two-dimensional: all surviving directions belong to a single line. A surviving nonzero lattice vector shows that this line has rational slope.

\begin{corollary}\label{cor:periodic-difference}
If $c$ is a nonperiodic finite-range rational configuration with a nonzero annihilator, then there are $x,y\in X_c$ such that $d=x-y$ is nonzero, vanishes on an open half-plane, and has a period parallel to its boundary.
\end{corollary}
\begin{proof}
Nonperiodicity makes all translates of $c$ distinct. Apply Lemma~\ref{lem:halfplane-pair}. By~\eqref{eq:inheritance}, the difference inherits an annihilator of $c$, so Proposition~\ref{prop:tangent-period} applies.
\end{proof}

\section{The annihilator ideal of the difference}\label{sec:ideal}
A period $qv$ gives an annihilator $\mon{qv}-1$. Theorem~\ref{thm:descent}, however, requires a generator of the \emph{whole} annihilator ideal. The vanishing half-plane supplies the additional information: it prevents a recurrence in the transverse direction from introducing a relation not already forced by the minimal polynomial in direction $v$.

\begin{theorem}\label{thm:exact-ideal}
Let $d:\Z^2\to\Q$ be nonzero. Suppose that $T^{qv}d=d$, where $v$ is primitive and $q\ge1$, and that $d$ vanishes on an open half-plane with boundary parallel to $v$. Then there is a monic nonconstant $\phi\in\Q[Z]$ such that
\begin{equation}\label{eq:exact-ideal}
 \phi\mid Z^q-1,
 \qquad \Ann(d)=\bigl(\phi(\mon v)\bigr).
\end{equation}
\end{theorem}
\begin{proof}
Complete $v$ to a lattice basis $v,w$, orient $w$ toward the other half-plane, and translate if necessary. In these coordinates we can write
\[
 d(a+q,b)=d(a,b),\qquad d(a,b)=0\quad(b<0).
\]
Translation leaves the annihilator unchanged. Let $Z,W$ denote the variables for shifts in the $v,w$ directions.

Choose a monic polynomial $\phi(Z)$ of least degree among the nonzero polynomials annihilating $d$ in the first coordinate. Division of $Z^q-1$ by $\phi$ leaves an annihilating remainder of smaller degree, which must vanish. Thus $\phi\mid Z^q-1$. Since $d\ne0$, this polynomial is nonconstant; since $Z^q-1$ is square-free over $\Q$, so is $\phi$, and $\phi(0)\ne0$.

We claim that every annihilator $F(Z,W)$ is divisible by $\phi$. After multiplication by a Laurent monomial, assume that $F$ is an ordinary polynomial. If some irreducible factor $\pi$ of $\phi$ does not divide $F$, put
\[
 e=(\phi/\pi)(T^v)d.
\]
This configuration is nonzero by minimality of $\phi$. It still vanishes on negative rows, and
\[
 \pi(T^v)e=0,\qquad F(T^v,T^w)e=0.
\]
The purpose of passing from $d$ to $e$ is that every polynomial in $Z$ not divisible by $\pi$ acts invertibly on the kernel of $\pi(T^v)$, by B\'ezout's identity.

Write $F(Z,W)=\sum_j f_j(Z)W^j$, let $j_*$ be the largest index with $\pi\nmid f_{j_*}$, and let $b_0$ be the first nonzero row of $e$. Evaluate $F(T^v,T^w)e=0$ on row $b_0-j_*$. Terms with $j>j_*$ vanish because their coefficients are multiples of $\pi$; terms with $j<j_*$ vanish because they sample rows below $b_0$. Thus the nonzero row $u=e(\,\cdot\,,b_0)$ satisfies
\[
 f_{j_*}(T^v)u=0,\qquad \pi(T^v)u=0.
\]
Choose $A,B\in\Q[Z]$ with $A\pi+Bf_{j_*}=1$. Applying this identity to $u$ gives $u=0$, a contradiction.

Every irreducible factor of $\phi$ therefore divides every coefficient of $F$ as a polynomial in $W$. Squarefreeness gives $\phi\mid F$. Restoring the Laurent monomial proves the same assertion in the Laurent ring. Conversely, every multiple of $\phi$ annihilates $d$. The change of lattice basis now yields~\eqref{eq:exact-ideal}.
\end{proof}

The minimal polynomial is necessary, not just a convenient normalization. For
$d(i,j)=\mathbf1_{\{j\ge0\}}$, the theorem gives $\Ann(d)=(X-1)$, although $X^2-1$ also annihilates $d$. On an $m\times n$ rectangle with $m\ge2$, the translated restrictions of $d$ span an $n$-dimensional space. The kernel of the map induced by $X^2-1$ has dimension $2n$. Choosing the larger period polynomial would therefore destroy the equality~\eqref{eq:full-kernel} on which the pattern count depends.

\begin{example}\label{ex:sharp}
The boundary term in Theorem~\ref{thm:descent} can be attained. Let
\[
 d(i,j)=\mathbf1_{\{j\ge i\}},\qquad \Phi=XY-1.
\]
The zero configuration belongs to $X_d$, and Theorem~\ref{thm:exact-ideal} gives $\Ann(d)=(XY-1)$. On $R_{m,n}$, the values of $j-i$ range over $m+n-1$ consecutive integers. Moving the threshold across these levels gives $m+n$ patterns, including the two constant ones. Since $\Phi(T)d=0$ and $|S|=(m-1)(n-1)$,
\[
 P_d(m,n)-P_{\Phi(T)d}(S)=m+n-1
 =mn-(m-1)(n-1)=|R|-|S|.
\]
\end{example}

\section{The case of two difference operators}\label{sec:twofactor}
We prove the result needed when the polynomial image of a minimal-area counterexample is periodic.

\begin{theorem}\label{thm:twofactor}
Let $c:\Z^2\to\Q$ have finite range. Suppose that
\[
 P_c(m,n)\le mn,\qquad \De_h\De_t c=0
\]
for some $m,n\ge1$ and nonzero $h,t\in\Z^2$. Then $c$ is periodic. If $h,t$ are nonparallel, a nonzero integer multiple of one of them is a period.
\end{theorem}

A sum of configurations with periods $h,t$ satisfies this mixed-difference identity. Szabados proved the corresponding two-periodic-summand theorem~\cite[Theorem~1]{Sz}. The argument below follows his proof: we retain the boundary selection, the ambiguous-extension count, and the finite description of strips, with their sources indicated at the relevant lemmas. The change is the periodicity extension in Lemma~\ref{lem:row-lifting}: the mixed-difference identity and a finite recurrence replace the general strip-extension result of~\cite[Lemma~3 and Appendix~A.2]{Sz}.

Here is the plan for nonparallel $h,t$. We find two translates agreeing on a sufficiently wide strip but differing on its next row, unless a multiple of $t$ is already a period. Low complexity supplies a local unique-extension rule on that boundary. The disagreement forces so few patterns inside the strip that Morse--Hedlund gives periodicity there. The local rule then extends periodicity to a few more rows. Only finitely many rows are needed: a difference of $c$ is $t$-periodic, and their common period consequently extends to the plane.

\subsection{One-dimensional lemmas}
\begin{lemma}\label{lem:finite-states}
Let $s:\Z\to A$ have finite range. Suppose that
\[
 s_i=s_j\quad\Longrightarrow\quad s_{i+1}=s_{j+1}
 \qquad(i,j\in\Z).
\]
Then $s$ has a positive period at most $|s(\Z)|$. The same conclusion holds with $i-1,j-1$ in place of $i+1,j+1$.
\end{lemma}
\begin{proof}
Successor defines a map on the finite set $s(\Z)$. It is surjective, since every occurrence has a predecessor, and hence is a permutation. The sequence follows one cycle of that permutation. Reversing the indices proves the other version.
\end{proof}

For completeness, we give the usual extension-count proof of the form of Morse--Hedlund used below; compare~\cite[Section~2.1]{CK}.

\begin{corollary}[Morse--Hedlund]\label{cor:morse}
If a finite-range bi-infinite word has at most $n$ words of length $n$, where $n\ge1$, then it has a positive period at most $n$.
\end{corollary}
\begin{proof}
Let $p(j)$ count its length-$j$ words, with $p(0)=1$. The counts are nondecreasing. If all $n$ increases up to $p(n)$ were strict, then $p(n)\ge n+1$. Hence $p(j+1)=p(j)$ for some $0\le j<n$. Restriction of occurring $(j+1)$-words to their first $j$ symbols is therefore bijective. Those $j$ symbols determine the following symbol, so the successive $(j+1)$-words satisfy the hypothesis of Lemma~\ref{lem:finite-states}. There are $p(j+1)\le n$ such words. Their period is a period of the original word.
\end{proof}

The next lemma is the finite-memory version with a periodic parameter. We state the unique-extension property only for inputs that occur; no rule on unused inputs is needed.

\begin{lemma}\label{lem:forcing}
Let $a:\Z\to A$ have finite range, and let $\eta:\Z\to E$ be periodic. Suppose that, for some $k\ge0$,
\[
 \eta_i=\eta_j,\quad a_{i+\ell}=a_{j+\ell}\ (0\le\ell<k)
 \quad\Longrightarrow\quad a_{i+k}=a_{j+k}.
\]
Then $a$ is periodic.
\end{lemma}
\begin{proof}
If $p\ge1$ is a period of $\eta$, the sequence
\[
 s_i=(i\bmod p,a_i,a_{i+1},\ldots,a_{i+k})
\]
has finite range and a uniquely determined successor. Indeed, equality of two such states gives equality of $\eta$ at the next indices and of the $k$ symbols needed there, and the hypothesis gives the one new symbol. Apply Lemma~\ref{lem:finite-states}. This also covers $k=0$.
\end{proof}

The output need not inherit a prescribed period of the parameter. For example, a constant parameter can accompany the recurrence $a_{i+1}=1-a_i$, whose nonconstant binary solutions have period two. We shall take common multiples of the finitely many row periods that arise, rather than try to preserve a prescribed period at each step.

\subsection{Unique extension across part of a boundary}
Fix a primitive $v$ and a complementary vector $w$, so that $v,w$ is a basis of $\Z^2$. The set $\{iv+jw:i\in\Z\}$ will be called row $j$. Points on a row are ordered by the integer $i$. These coordinates describe subsets of the original lattice; a rectangle in the original coordinates need not be a rectangle in $(i,j)$ coordinates.

The following is a variant of the edge-removal construction in~\cite[Lemma~2]{Sz}, followed by the unique-extension step from the proof of~\cite[Lemma~4]{Sz}. We delete a shorter extreme row at each stage. The interpolation argument of~\cite[Lemma~2.11]{CK} then bounds the lengths of the remaining rows, as needed for Morse--Hedlund.

\begin{lemma}\label{lem:boundary-window}
Suppose $P_c(R)\le|R|$ for a nonempty rectangle $R$. After choosing the sign of $w$ and translating a subset of $R$, there is a finite set $B$ with bottom row
\[
 E=\{0,v,\ldots,(e-1)v\},\qquad e\ge1,
\]
and top row $H\ge0$ such that, with $C=B\setminus E$,
\begin{equation}\label{eq:boundary-cost}
 P_c(B)\le|B|,\qquad r:=P_c(B)-P_c(C)<e.
\end{equation}
If $H>0$, each integer row from $0$ to $H$ contains at least $e-1$ consecutive points of $B$. Furthermore, for some $0\le k<e$, put
\begin{equation}\label{eq:boundary-rule-domain}
 D_k=C\cup\{0,v,\ldots,(k-1)v\},
\end{equation}
where the second set is empty for $k=0$. Then, for every $u,u'\in\Z^2$,
\begin{equation}\label{eq:boundary-rule}
 \restr{T^u c}{D_k}=\restr{T^{u'}c}{D_k}
 \quad\Longrightarrow\quad c(u+kv)=c(u'+kv).
\end{equation}
\end{lemma}
\begin{proof}
Starting from $R$, repeatedly delete the nonempty extreme row with fewer sites, choosing either in a tie. Each remaining set is the set of lattice points in the real rectangle intersected with a closed strip. Its nonempty rows are segments of consecutive lattice sites.

Suppose that the current extreme rows have indices $j_-<j_+$ and tangential intervals $[\alpha_-,\beta_-]$ and $[\alpha_+,\beta_+]$. If the shorter row has $e$ sites, both interval lengths are at least $e-1$. For $j_-<j<j_+$, let $\lambda=(j-j_-)/(j_+-j_-)$. Convexity puts the interval
\[
 [(1-\lambda)\alpha_-+\lambda\alpha_+,\,
   (1-\lambda)\beta_-+\lambda\beta_+]
\]
in row $j$ of the real set. Its length is at least $e-1$, so it contains at least $e-1$ consecutive integers. They give points of the current window, proving the row-length bound at each deletion step.

The discrepancy begins at $\delta_c(R)\le0$ and ends at $\delta_c(\varnothing)=1$. At the first step from nonpositive to positive discrepancy, write the deletion as $B\to C=B\setminus E$. Then
\[
 0<\delta_c(C)-\delta_c(B)
 =e-\bigl(P_c(B)-P_c(C)\bigr),
\]
which proves~\eqref{eq:boundary-cost}. Reverse $w$ if necessary to place the deleted row below the others, and translate its first point to the origin.

For $0\le i\le e$, let $D_i=C\cup\{0,v,\ldots,(i-1)v\}$. The $e$ nonnegative integer increments $P_c(D_{i+1})-P_c(D_i)$ sum to $r<e$. Some increment is zero. At that index $k$, restriction from occurring $D_{k+1}$-patterns to occurring $D_k$-patterns is a bijection. This is~\eqref{eq:boundary-rule}.
\end{proof}

The row-length bound is $e-1$, not $e$: the endpoints of the interpolated interval need not be integers. In the $(v,w)$ coordinates, the set
\[
 \{(0,0),(1,0),(1,1),(1,2),(2,2)\}
\]
illustrates the loss. Its extreme rows have two sites each, but its middle row has only one.

\subsection{Two translates agreeing on a strip}
Assume in the next lemmas that
\begin{equation}\label{eq:two-normalized}
 h=qv,\qquad t=av+Mw,\qquad q,M\ge1,\quad a\in\Z,
 \qquad \De_h\De_t c=0.
\end{equation}
The identity says that $\De_h c$ is $t$-periodic. Thus the differences between the translates $T^{jt}c$ are $h$-periodic. As in Szabados's proof~\cite[Section~6]{Sz}, a finite block consequently distinguishes their restrictions to an infinite strip. We express the resulting alternative using Lemma~\ref{lem:finite-states}.

\begin{lemma}\label{lem:strip-states}
For every $W\ge M$, either a nonzero integer multiple of $t$ is a period of $c$, or there are translates $x,y$ of $c$ with
\begin{equation}\label{eq:strip-pair}
 \De_h(x-y)=0,\qquad x=y\text{ on rows }1,\ldots,W,
 \qquad x\ne y\text{ somewhere on row }0.
\end{equation}
\end{lemma}
\begin{proof}
Let $S_W$ be the union of rows $1,\ldots,W$, and put
\[
 \sigma_j=\restr{T^{jt}c}{S_W},\qquad
 F_W=\{iv+\ell w:0\le i<q,\ 1\le\ell\le W\}.
\]
For any $j,k$, the difference $T^{jt}c-T^{kt}c$ is $qv$-periodic. Its vanishing on $F_W$ is therefore equivalent to its vanishing on $S_W$. Restriction to $F_W$ is thus injective on the set $\{\sigma_j:j\in\Z\}$, which is consequently finite. The individual $\sigma_j$ need not be $qv$-periodic.

If $\sigma_j=\sigma_k$ always implies $\sigma_{j-1}=\sigma_{k-1}$, Lemma~\ref{lem:finite-states} gives a positive period $p$ of $(\sigma_j)$. Since $W\ge M$, the strips $S_W+jt$ cover $\Z^2$. Equality of the strip restrictions at indices $j$ and $j+p$ then gives $T^{pt}c=c$ everywhere.

Otherwise, choose $j,k$ with $\sigma_j=\sigma_k$ but $\sigma_{j-1}\ne\sigma_{k-1}$. Set $x_0=T^{jt}c$ and $y_0=T^{kt}c$. They agree on $S_W$ but differ on $S_W-t$, whose rows range from $1-M$ to $W-M$. Every positive row in that range already agrees. Let $s\in\{1-M,\ldots,0\}$ be the largest index of a differing row in this range, and put
\[
 x=T^{sw}x_0,\qquad y=T^{sw}y_0.
\]
They differ on row zero. The old rows $s+1,\ldots,0$ agree by the choice of $s$, and the old rows $1,\ldots,W$ agree by assumption. Since $s\le0$, this includes the old rows $s+1,\ldots,s+W$, giving the required agreement on the new rows $1,\ldots,W$. Their difference remains $h$-periodic.
\end{proof}

\subsection{Counting extensions and applying Morse--Hedlund}
We adapt the argument of Cyr and Kra~\cite[Lemma~2.24]{CK}, also used by Szabados~\cite[Lemma~4]{Sz}. Grouping finitely many rows into one word converts the extension bound into a word-complexity bound. Here the periodic difference permits a single forward boundary rule. Fix $B$ and the notation of Lemma~\ref{lem:boundary-window}, with $H\ge1$.

\begin{lemma}\label{lem:periodic-interior}
Suppose that translates $x,y$ of $c$ agree on rows $1,\ldots,H$, differ on row zero, and have $qv$-periodic difference for some $q\ge1$. Then $r\ge1$, and $x$ restricted to rows $1,\ldots,H$ has a common period $pv$ with $1\le p\le r$.
\end{lemma}
\begin{proof}
The two configurations cannot agree on $E+iv$ for any $i\in\Z$. Such agreement and~\eqref{eq:boundary-rule} would extend their boundary agreement indefinitely to the right. To obtain the next point, place $kv$ there: the preceding $k<e$ boundary values agree, and all values on the corresponding translate of $C$ agree. For $k=0$, no preceding boundary values are needed. Since the boundary difference is $qv$-periodic, vanishing on a right half-line would imply vanishing on the whole boundary, contrary to assumption.

Thus every $C$-pattern of $x$ encountered under shifts by $iv$ has at least two distinct $B$-extensions in $c$. Let $\pi:\Pat_c(B)\to\Pat_c(C)$ be restriction. The number of patterns with two or more extensions is at most
\[
 \sum_{p\in\Pat_c(C)}\bigl(|\pi^{-1}(p)|-1\bigr)
 =P_c(B)-P_c(C)=r.
\]
At least one such pattern occurs, so $r\ge1$. Since $r<e$, each row of $C$ has at least $e-1\ge r$ consecutive sites.

To obtain one period for all these rows, regard them together as a word over a product alphabet. If $\alpha_jv+jw$ is the first site of row $j$ of $C$, define
\[
 b_i=\bigl(x((\alpha_j+i)v+jw)\bigr)_{j=1}^{H},\qquad i\in\Z.
\]
Every length-$r$ word of $b$ is a restriction of one of the translated $C$-patterns just counted. There are at most $r$ such words. Corollary~\ref{cor:morse} gives a period $p\le r$ for $b$, and each coordinate gives the period $pv$ on its entire row.
\end{proof}

The argument uses only a forward unique-extension rule. Agreement propagates to a half-line, and the periodic difference then extends it to the entire boundary. This is the point at which the mixed-difference hypothesis enters the boundary count.

\subsection{Extending row periods}
\begin{lemma}\label{lem:row-lifting}
Assume~\eqref{eq:two-normalized} and the unique-extension property~\eqref{eq:boundary-rule}, with $H\ge1$. If rows $1,\ldots,H$ of $c$ are individually periodic in direction $v$, then a nonzero integer multiple of $h$ is a period of $c$.
\end{lemma}
\begin{proof}
Suppose that the $H$ rows immediately above row $j$ are periodic. Set
\[
 a_i=c(iv+jw),\qquad \eta_i=\restr{T^{iv+jw}c}{C}.
\]
A common positive period of those $H$ rows makes $\eta$ periodic. Property~\eqref{eq:boundary-rule} says exactly that $\eta_i$ and $a_i,\ldots,a_{i+k-1}$ determine $a_{i+k}$. Lemma~\ref{lem:forcing} shows that row $j$ is periodic as well.

Starting from rows $1,\ldots,H$, repeat this argument downward until $M$ consecutive rows are periodic. This requires at most $\max(0,M-H)$ new rows. Choose a positive integer $p$ divisible by $q$ and by a period of each of those $M$ rows, and set $d=\De_{pv}c$. The configuration $d$ vanishes on the chosen rows. Writing $p=Lq$, we have
\begin{equation}\label{eq:telescoping-lift}
 d=\sum_{i=0}^{L-1}T^{iqv}\De_{qv}c.
\end{equation}
Every summand is $t$-periodic by~\eqref{eq:two-normalized}, so $d$ is $t$-periodic. Translation by $t$ changes the row index by $M$. Every site can therefore be moved along its $t$-orbit to one of the $M$ consecutive rows on which $d$ vanishes. Hence $d=0$ everywhere, and $pv=Lh$ is a nonzero period.
\end{proof}

Only finitely many row periods are combined. Individual periodicity of every row would not by itself imply a common period: for example,
\[
 c(i,j)=\mathbf1_{\{i\equiv0\pmod{|j|+2}\}}
\]
has no nonzero horizontal period. In the lemma, it is the preexisting $t$-periodicity of $\De_{qv}c$ that propagates a common period from finitely many rows.

\begin{proof}[Proof of Theorem~\ref{thm:twofactor}]
If $h,t$ are parallel, reverse one vector if necessary and take a common positive integer multiple $u=ah=bt$, with $a,b\ge1$. The telescoping expressions for $\De_u$ in terms of $\De_h$ and $\De_t$ give $\De_u^2c=0$. Lemma~\ref{lem:repeated} yields $\De_uc=0$.

Suppose now that $h,t$ are nonparallel. Write $h=qv$ with $v$ primitive and $q\ge1$, and complete $v$ to a lattice basis. Apply Lemma~\ref{lem:boundary-window} to the witnessing rectangle to obtain $B$ and a sign of $w$. Reverse $t$ if necessary so that $t=av+Mw$ with $M>0$. This preserves the identity because $\De_{-t}=-T^{-t}\De_t$.

If $H=0$, then $B=E$ and $P_c(E)\le e$. Every row word has at most $e$ words of length $e$. By Corollary~\ref{cor:morse}, each row has a period at most $e$, so $q(e!)v$ is a period of the whole configuration and a multiple of $h$.

For $H\ge1$, choose $W\ge\max(H,M)$. Lemma~\ref{lem:strip-states} either gives a period that is a multiple of $t$, or gives translates $x,y$ satisfying~\eqref{eq:strip-pair}. In the latter case, Lemma~\ref{lem:periodic-interior} makes rows $1,\ldots,H$ of $x$ periodic. The mixed identity and unique-extension property hold for $x$, so Lemma~\ref{lem:row-lifting} gives it a period that is a multiple of $h$. A period of a translate is a period of $c$.
\end{proof}

\section{Proof of Nivat's conjecture}\label{sec:completion}
\begin{proof}[Proof of Theorem~\ref{thm:main}]
Suppose a counterexample exists. After rationally labeling its alphabet, choose a nonperiodic finite-range rational configuration $c$ with a witnessing rectangle $R_{m,n}$ of least possible area $N=mn$. The minimum is taken over all finite rational alphabets.

Lemma~\ref{lem:ann-exists} gives a nonzero annihilator of $c$. By Corollary~\ref{cor:periodic-difference}, there are $x,y\in X_c$ whose difference $d$ is nonzero, vanishes on a half-plane, and has a period $qv$ parallel to its boundary, with $v$ primitive. Theorem~\ref{thm:exact-ideal} gives
\[
 \Ann(d)=(\Phi),\qquad
 \Phi=\phi(\mon v),\qquad \phi\mid Z^q-1,
\]
where $\phi$ is monic, nonconstant, and has nonzero constant term.

The configuration $c'=\Phi(T)c$ still has finite rational range. Corollary~\ref{cor:line-descent} gives a nonempty rectangle $S$ of area less than $N$ with $P_{c'}(S)\le|S|$. Translating $S$ if needed, minimality of $N$ makes $c'$ periodic. Choose a nonzero period $t$ of $c'$. Then
\[
 \De_t\phi(T^v)c=0.
\]
Multiply by $\psi(T^v)$, where $\psi(Z)=(Z^q-1)/\phi(Z)$. The result is
\[
 \De_t\De_{qv}c=0.
\]
Theorem~\ref{thm:twofactor}, applied to $c$ and its original rectangular complexity bound, makes $c$ periodic, a contradiction.
\end{proof}

The two uses of the alphabet and rectangle in this argument are different. The alphabet may change under $\Phi(T)$, which is why the minimum ranges over all finite rational alphabets. The smaller rectangle remains in the original lattice coordinates, so no complexity assertion on a period sublattice is required.

\appendix
\section{A proof of the product-annihilator statement}\label{app:product}
We prove Proposition~\ref{prop:product}. The arithmetic step is Lemma~7 of Kari and Szabados~\cite{KS}, and the resulting product is the one in their Theorem~11. Their Lemma~8 uses the dilation relations in a Vandermonde argument at common zeros. The identity below carries out the corresponding elimination directly in the Laurent-polynomial ring.

\begin{proof}[Proof of Proposition~\ref{prop:product}]
Scale $c$ and a nonzero annihilator separately so that both are integral, and write
\[
 f=\sum_{i=0}^{k-1}a_i\mon{u_i}
\]
with distinct $u_i$ and nonzero integers $a_i$. Since $c\ne0$, a monomial cannot annihilate it, so $k\ge2$. Let
\[
 K=\max_z|c(z)|\sum_i|a_i|,\qquad \rho=K!,\qquad
 f_n=\sum_i a_i\mon{nu_i}\quad(n\ge1).
\]
The positive integer $K$ bounds every value of $f_n(T)c$, independently of $n$.

If $f_n$ annihilates $c$ and $p>K$ is prime, reduction modulo $p$ gives
\[
 f_{np}\equiv f_n^{\,p}\pmod p.
\]
Thus every value of $f_{np}(T)c$ is divisible by $p$ and has absolute value at most $K<p$. It vanishes. Starting with $f_1=f$ and factoring an integer into primes shows that $f_n$ annihilates $c$ whenever $\gcd(n,\rho)=1$. In particular,
\begin{equation}\label{eq:dilation-family}
 f_{1+j\rho}\in\Ann(c)\qquad(j\ge0).
\end{equation}

Put $M_i=\mon{\rho u_i}$. The relations in~\eqref{eq:dilation-family} have the form
$f_{1+j\rho}=\sum_i a_i\mon{u_i}M_i^j$. We can therefore take linear combinations of them using the coefficients of a polynomial that vanishes at $M_1,\ldots,M_{k-1}$. Namely, write
\[
 H(U)=\prod_{i=1}^{k-1}(U-M_i)=\sum_{j=0}^{k-1}H_jU^j,
 \qquad H_j\in\LL.
\]
Then
\begin{align}
 \sum_{j=0}^{k-1}H_jf_{1+j\rho}
 &=\sum_{i=0}^{k-1}a_i\mon{u_i}H(M_i)\notag\\
 &=a_0\mon{u_0}\prod_{i=1}^{k-1}(M_0-M_i).
 \label{eq:interpolation}
\end{align}
The left side annihilates $c$. The scalar and monomial factors on the right are units in $\LL$. Removing them gives the annihilator
\[
 \prod_{i=1}^{k-1}\bigl(\mon{\rho(u_i-u_0)}-1\bigr).
\]
All vectors $\rho(u_i-u_0)$ are nonzero, as required.
\end{proof}

For a trinomial $f=\alpha+\beta X+\gamma Y$, identity~\eqref{eq:interpolation} reads
\[
 f_{1+2\rho}-(X^\rho+Y^\rho)f_{1+\rho}+X^\rho Y^\rho f_1
 =\alpha(1-X^\rho)(1-Y^\rho).
\]
No separation of the directions of the resulting binomials is needed in Proposition~\ref{prop:tangent-period}; boundedness removes repetitions after the transverse factors have been cancelled.

\subsection*{Acknowledgement}
An AI language model was used in developing the arguments, drafting the text, and reviewing the exposition and references.

\begin{thebibliography}{ELMW}
\bibitem[BL]{BL}
M.~Boyle and D.~Lind,
\emph{Expansive subdynamics},
Trans. Amer. Math. Soc. \textbf{349} (1997), no.~1, 55--102.
\href{https://doi.org/10.1090/S0002-9947-97-01634-6}{doi:10.1090/S0002-9947-97-01634-6}.

\bibitem[CK]{CK}
V.~Cyr and B.~Kra,
\emph{Nonexpansive $\mathbb Z^2$-subdynamics and Nivat's conjecture},
Trans. Amer. Math. Soc. \textbf{367} (2015), no.~9, 6487--6537.
\href{https://doi.org/10.1090/S0002-9947-2015-06391-0}{doi:10.1090/S0002-9947-2015-06391-0}.
Numbering refers to \href{https://arxiv.org/abs/1208.4090v2}{arXiv:1208.4090v2} (2013).

\bibitem[ELMW]{ELMW}
M.~Einsiedler, D.~Lind, R.~Miles, and T.~Ward,
\emph{Expansive subdynamics for algebraic $\mathbb Z^d$-actions},
Ergodic Theory Dynam. Systems \textbf{21} (2001), no.~6, 1695--1729.
\href{https://doi.org/10.1017/S014338570100181X}{doi:10.1017/S014338570100181X}.

\bibitem[KM]{KM}
J.~Kari and E.~Moutot,
\emph{Decidability and periodicity of low complexity tilings},
Theory Comput. Syst. \textbf{67} (2023), 125--148.
\href{https://doi.org/10.1007/s00224-021-10063-8}{doi:10.1007/s00224-021-10063-8}.

\bibitem[KS]{KS}
J.~Kari and M.~Szabados,
\emph{An algebraic geometric approach to Nivat's conjecture},
Inform. and Comput. \textbf{271} (2020), article 104481.
\href{https://doi.org/10.1016/j.ic.2019.104481}{doi:10.1016/j.ic.2019.104481}.
Numbered statements cited here refer to
\href{https://arxiv.org/abs/1605.05929v1}{arXiv:1605.05929v1} (2016).

\bibitem[MH]{MH}
M.~Morse and G.~A.~Hedlund,
\emph{Symbolic dynamics},
Amer. J. Math. \textbf{60} (1938), no.~4, 815--866.
\href{https://doi.org/10.2307/2371264}{doi:10.2307/2371264}.

\bibitem[Sz]{Sz}
M.~Szabados,
\emph{Nivat's conjecture holds for sums of two periodic configurations},
in \emph{SOFSEM 2018: Theory and Practice of Computer Science},
Lecture Notes in Comput. Sci., vol.~10706, Springer, Cham, 2018, pp.~539--551.
Numbered statements cited here refer to
\href{https://arxiv.org/abs/1710.05360v1}{arXiv:1710.05360v1} (2017).
\end{thebibliography}
\end{document}
